\section{Interior lattice points and interpolation}
\label{sec:saturation}

The reshaping proposition preserves the number of monomials, but does not
describe each resulting exponent. For interpolation we need more than the
correct asymptotic number: we need specific sets of exponents. Additivity
across degrees supplies the missing information. An additive family $F_k$
occupying asymptotically the full volume of $kP$ contains every lattice point
of $kK$ for each fixed compact set $K\subset\operatorname{int}P$, once the
degree $k$ is sufficiently large.

The general semigroup approximation theorem of
\citet[Theorem~1.6]{KavehKhovanskii2012} provides a broader setting for
interior-saturation statements. Under the full-density hypothesis needed
here, the following direct proof uses only decompositions into two summands.

\begin{lemma}[Interior saturation]
\label{lem:interior-saturation}
Let $n\ge1$ be an integer, let $P\subset\R^n$ be a compact convex set with
nonempty interior, and let $F_k\subset kP\cap\Z^n$ for every integer
$k\ge1$. Suppose that
\[
 F_k+F_l\subset F_{k+l}\quad(k,l\ge1),
 \qquad
 \#F_k=\vol(P)k^n+o(k^n).
\]
For every compact set $K\subset\operatorname{int}P$, there is an integer
$k_K$ such that
\[
 kK\cap\Z^n\subset F_k\qquad(k\ge k_K).
\]
\end{lemma}

\begin{proof}
The lattice-point estimate for a full-dimensional compact convex set gives
\[
 \#(kP\cap\Z^n)=\vol(P)k^n+o(k^n).
\]
Consequently the sets of missing exponents
\[
 H_k=(kP\cap\Z^n)\setminus F_k
\]
satisfy $\#H_k=o(k^n)$. This estimate alone would allow individual missing
points deep inside $kP$. We use additivity to exclude them.

We may assume that $K$ is nonempty. Choose $\delta>0$ so that every closed
ball of radius $\delta$ centered in $K$ is contained in $P$. For a large
integer $k$, put $k_1=\lfloor k/2\rfloor$ and $k_2=k-k_1$, and fix
$y\in kK\cap\Z^n$. Write $z=y/k$. Every integer vector $y_1$ satisfying
\[
 \lVert y_1-k_1z\rVert\le\delta\min(k_1,k_2)
\]
gives a decomposition $y=y_1+y_2$ with $y_i\in k_iP\cap\Z^n$. Indeed,
both $y_1/k_1$ and $y_2/k_2$ are within distance $\delta$ of $z$.
The ball in this display contains at least $c k^n$ lattice points for all
sufficiently large $k$, where $c>0$ is independent of $y$.

At most $\#H_{k_1}$ choices have $y_1\notin F_{k_1}$, and at most
$\#H_{k_2}$ choices have $y_2\notin F_{k_2}$, because $y_1\mapsto y-y_1$
is injective. The total number excluded is $o(k^n)$. Thus some decomposition
has $y_i\in F_{k_i}$, and additivity gives $y\in F_k$. All constants and
degree thresholds used here are uniform over $y\in kK\cap\Z^n$.
\end{proof}

The lemma requires no finite-generation hypothesis on the sets $F_k$ or on
the semigroup they form. It also permits real, rather than rational,
bounding polytopes. We next describe the exponents needed for interpolation
in the simplex that will result from repeated reshaping.

\begin{lemma}[Interpolation in an elongated simplex]
\label{lem:simplex-interpolation}
Let $n\ge2$, $r\ge1$, and $m\ge0$ be integers, and let
$c_1,\ldots,c_r\in\C^*$ be distinct. Put
\[
 E_{r,m}=
 \left\{(\beta,b)\in\Z_{\ge0}^{n-1}\times\Z_{\ge0}:
       |\beta|+\frac br<m+1\right\}.
\]
Then $\M(E_{r,m})$ surjects onto jets through degree $m$ at the points
\[
 q_i=(1,\ldots,1,c_i)\in\Torus^n,\qquad 1\le i\le r.
\]
\end{lemma}

\begin{proof}
Write the coordinates as $(x',Y)$, with $n-1$ coordinates in $x'$. At $q_i$,
an arbitrary jet through degree $m$ has a unique expression
\[
 \sum_{|\alpha|\le m}(x'-1)^\alpha h_{i,\alpha}(Y-c_i),
 \qquad
 \deg h_{i,\alpha}\le m-|\alpha|.
\]
For each $\alpha$, the Chinese remainder theorem provides a polynomial
$f_\alpha(Y)$ such that
\[
 \begin{split}
 f_\alpha(Y)&\equiv h_{i,\alpha}(Y-c_i)
       \pmod{(Y-c_i)^{m+1-|\alpha|}}\quad(1\le i\le r),\\
 \deg f_\alpha&<r(m+1-|\alpha|).
 \end{split}
\]
The polynomial
\[
 f(x',Y)=\sum_{|\alpha|\le m}(x'-1)^\alpha f_\alpha(Y)
\]
has all the prescribed jets. Every monomial $(x')^\beta Y^b$ occurring in
its $\alpha$-summand satisfies $\beta\le\alpha$ componentwise and hence
\[
 |\beta|+\frac br
 \le |\alpha|+\frac br
 <|\alpha|+(m+1-|\alpha|)=m+1.
\]
Thus $f\in\M(E_{r,m})$.
\end{proof}

The two lemmas now convert asymptotically full additive supports into
simultaneous jet surjectivity. Translating a slightly smaller simplex into
the interior avoids any assertion about exponents on the boundary.

\begin{proposition}[Jets from full-volume supports]
\label{prop:asymptotic-jets}
Let $n\ge2$ and $r\ge1$ be integers, let $w>0$, and set
\[
 P=\Delta(w,\ldots,w,rw).
\]
Suppose that finite sets $F_k\subset kP\cap\Z^n$, indexed by all integers
$k\ge1$, satisfy
\[
 F_k+F_l\subset F_{k+l},
 \qquad
 \#F_k=\vol(P)k^n+o(k^n).
\]
For every $0<\theta<w$ and every choice of distinct
$c_1,\ldots,c_r\in\C^*$, the space $\M(F_k)$ surjects onto jets through
degree $\lfloor k\theta\rfloor$ at $(1,\ldots,1,c_i)$ for all sufficiently
large $k$.
\end{proposition}

\begin{proof}
Choose $\theta/w<\lambda<1$. There is a vector $a$ with strictly positive
coordinates such that $a+\lambda P\subset\operatorname{int}P$: it suffices
to require
\[
 \frac{a_1+\cdots+a_{n-1}}w+\frac{a_n}{rw}<1-\lambda.
\]
Choose $a_k\in k^{-1}\Z^n$ with $a_k\to a$. For all sufficiently large
$k$, the translated simplices $a_k+\lambda P$ lie in a single compact
subset $K$ of $\operatorname{int}P$. Lemma~\ref{lem:interior-saturation}
therefore gives
\[
 ka_k+(k\lambda P\cap\Z^n)\subset F_k.
\]
Consequently $\M(F_k)$ contains the monomial space for
$k\lambda P\cap\Z^n$, multiplied by the monomial with exponent $ka_k$.
Multiplication by this monomial is an isomorphism on each jet algebra at
a torus point, since the monomial is a unit there.

Put $m=\lfloor k\theta\rfloor$. For all sufficiently large $k$,
$m+1\le k\lambda w$, and hence
\[
 E_{r,m}\subset k\lambda P\cap\Z^n.
\]
Lemma~\ref{lem:simplex-interpolation} supplies jet surjectivity for this
subspace. The monomial translation preserves surjectivity and proves the
claim.
\end{proof}
